\section{Plan de plataforma, seguridad y componentes comunes}
La distribución inicial de las 2.024 HH de estas cuentas incorpora aprobación de arquitectura, seguridad y modelo en H2, preparación de cinco ambientes y observabilidad para H3, construcción de componentes comunes E1 desde el mes 8 y adaptadores administrativos para E2. Los períodos de las fichas son ventanas de ejecución propuestas; las esperas del artículo 18 se incorporarán a sus cadenas de recepción y no se consideran absorbidas por las horas de configuración \parencite[arts.~17--18; Formulario~E-25; pp.~12--74]{bases-admin}.

Las actividades del registro \path{componentes/actividades_plataforma.csv} corresponden a los 150 aportes por perfil de las 41 fichas. Cada actividad conserva ejecutor, responsable último, cuenta, lugar, ventana y predecesoras. El aporte de Calidad consolida la comprobación de los aportes del paquete; las tareas técnicas dependen de los insumos y contratos de las fichas previas. Su suma por paquete y perfil reproduce las mismas horas, sin porcentaje adicional de gestión ni duplicación del componente padre. La coincidencia de dos fichas en una quincena no acredita que sus actividades secuenciales quepan en ella: la red deberá resolver dedicación, duración y recepción.

Los paquetes de configuración incluyen revisión y comprobación técnica de su objeto, sin ejecutar todas las pruebas del sistema. Desarrollo incluye implementación acotada, pruebas unitarias y documentación de componentes comunes; las pantallas y reglas de cada dominio permanecen en sus cuentas. Las habilitaciones de terceros se registran como condiciones de inicio, y las visitas o configuración en recinto respetarán congelamiento, campañas, fechas náuticas y restricciones marítimas. El enrolamiento se agrega por lotes del inventario efectivo, sin atribuir al cliente las horas técnicas de Synaptix.

\begingroup\small\setlength{\tabcolsep}{2pt}
\begin{longtable}{R{9mm}R{12mm}R{12mm}R{12mm}R{12mm}R{12mm}R{12mm}R{14mm}R{12mm}R{14mm}}
\caption{Cargas propuestas de plataforma, seguridad y servicios comunes}\label{tab:sd7-plataforma-curva}\\\hline
Mes & A & D & S & Q & V & I & O+L & J & HH\\\hline\endfirsthead
\hline Mes & A & D & S & Q & V & I & O+L & J & HH\\\hline\endhead
3 & 0 & 0 & 8 & 0 & 0 & 24 & 0 & 8 & 40 \\\hline
4 & 24 & 0 & 72 & 8 & 0 & 80 & 24 & 8 & 216 \\\hline
5 & 0 & 24 & 88 & 48 & 0 & 112 & 112 & 0 & 384 \\\hline
6 & 8 & 32 & 104 & 96 & 0 & 80 & 160 & 8 & 488 \\\hline
8 & 0 & 16 & 56 & 64 & 296 & 48 & 0 & 0 & 480 \\\hline
9 & 8 & 16 & 40 & 40 & 168 & 32 & 0 & 0 & 304 \\\hline
15 & 0 & 8 & 8 & 16 & 64 & 16 & 0 & 0 & 112 \\\hline
\end{longtable}\FuenteTabla{Suma por perfil de las fichas, sin horas de habilitaciones del cliente ni paquetes variables de enrolamiento. Las ventanas requieren nivelación completa.}\endgroup


\subsection{Balance parcial de las cuentas cuantificadas}
La Tabla~\ref{tab:sd7-balance-conocido} suma plataforma, seguridad y servicios comunes con las cuentas cuantificadas de migración, ensayos/implantación e innovación, incluyendo seguimiento fijo. Sustituye, para ese alcance conjunto, el contraste de innovación aislada. La hipótesis de 144 HH productivas por persona y mes se conserva como referencia; las funciones I y O+L se suman por titular.
\begin{longtable}{R{18mm}L{36mm}R{30mm}R{30mm}}
\caption{Sobrecargas de las cuentas cuantificadas}\label{tab:sd7-balance-conocido}\\\hline
 \EncabezadoTabla{Mes} & \EncabezadoTabla{Perfil/persona} & \EncabezadoTabla{HH conocidas} & \EncabezadoTabla{Exceso sobre 144}\\\hline\endfirsthead
\hline \EncabezadoTabla{Mes} & \EncabezadoTabla{Perfil/persona} & \EncabezadoTabla{HH conocidas} & \EncabezadoTabla{Exceso sobre 144}\\\hline\endhead
6 & O+L & 160 & 16 \\\hline
8 & D & 160 & 16\\\hline
8 & V & 344 & 200 \\\hline
9 & V & 264 & 120\\\hline
16 & V & 272 & 128 \\\hline
17 & Q & 156 & 12\\\hline
17 & V & 384 & 240 \\\hline
18 & Q & 168 & 24\\\hline
20 & Q & 152 & 8 \\\hline
21 & O+L & 164 & 20\\\hline
\end{longtable}\FuenteTabla{Plataforma/seguridad/comunes, migración, ensayos/implantación, innovaciones con mes fijo y seguimiento fijo. Excluye trabajos variables y demás cuentas. Una persona por titular de I y O+L.}


En el mes 6, la configuración y comprobación concentra 160 HH en Operación; en el 8, Datos suma 160. Desarrollo suma 344 en el 8 y 264 en el 9, además de sus excesos de 16--17. Este contraste parcial excluye núcleos funcionales, capacitación general y pruebas integrales; la asignación conjunta posterior incorpora sus paquetes fijos y propone refuerzos. La suficiencia total requiere además calendario, puertas, trabajos variables y cobertura. La construcción desde el mes 8, comprometida en SD1, impide resolver el exceso trasladando indiscriminadamente programación al 6 o 7. La Dirección deberá asegurar la provisión y el costo de los refuerzos propuestos, comprobar competencias y disponibilidad e incorporar cargas variables antes de fijar la línea base. No se supone que un perfil de seguridad, datos o arquitectura cubra automáticamente programación de todos los módulos.

Los 2.024 HH de estas cuentas se suman a las 4.836 de implementación ya cuantificadas: subtotal parcial 6.860 HH. Al incorporar núcleos funcionales y gobierno/pruebas/formación, la implementación fija conjunta alcanza 13.588 HH, más enrolamiento de terminales, revisión documental y otros trabajos variables. Este subtotal no constituye el precio ni esfuerzo del contrato completo y no absorbe operación. Las cargas de servicio de 21--56 conservarán su registro independiente, incluido mantenimiento de estos componentes.
